\subsection*{§. 201.\\ 椭圆几何与双曲几何中的 Frenet 公式。}

本章最后的研究，将把第一章所述的曲线理论推广到三维椭圆度量与双曲度量。我们只建立这些空间中的\emph{Frenet 公式}\srcfn{1}{这些公式最初见于我在上文引用的《年刊》论文。}；借助这些基本公式，读者可以重新考察第一章的各项研究，并观察空间曲率对最后结果所造成的轻微改变。

考察三维椭圆空间中的一条曲线 \(C\)，以其动点的 Weierstrass 坐标
\[
 x_0,x_1,x_2,x_3
\]
作为弧长 \(s\) 的函数来规定它。若以 \(\alpha_0,\alpha_1,\alpha_2,\alpha_3\) 表示曲线在 \((x_i)\) 处的法平面的坐标，则
\[
\tag{45}
 \alpha_i=R\frac{dx_i}{ds},\qquad (i=0,1,2,3);
\]
\(\alpha_i\) 也称为切线的方向余弦。先来计算曲率 \(\frac1\rho\)，即曲线 \(C\) 的曲率。采用 §. 166 中 \(\frac1\rho\) 的一般定义：它是比值 \(\frac{2d}{t^2}\) 的极限，其中 \(d\) 表示两点 \(M',\overline{M}\) 之间的距离；这两点分别由在曲线上和在 \(M\) 点的切线上截取无穷小弧 \(=t\) 而得。我们有
\[
 x'_i=x_i+\frac tR\alpha_i+\frac{t^2}{2R}\frac{d\alpha_i}{ds}+\cdots
\]
\[
 \overline{x}_i=x_i\cos\left(\frac tR\right)+\alpha_i\sen\left(\frac tR\right)
\]
或
\[
 \overline{x}_i=x_i+\frac tR\alpha_i-\frac{t^2}{2R^2}x_i+\cdots,
\]
这里略去了 \(t\) 的二次以上各次幂。于是
\[
 x'_i-\overline{x}_i=\frac{t^2}{2R}\left(\frac{d\alpha_i}{ds}+\frac{x_i}{R}\right)+\cdots,
 \qquad
 d=R\sqrt{\mathop{\mathrm{S}}(x'_i-\overline{x}_i)^2},
\]
从而
\[
 \frac1\rho=\lim\frac{2d}{t^2}=\sqrt{\mathop{\mathrm{S}}\left(\frac{d\alpha_i}{ds}+\frac{x_i}{R}\right)^2}.
\]
这里同欧几里得几何中一样，我们约定只赋予曲率一个绝对值，因而取根式的正值。

同时可见，曲线主法线的方向余弦 \(\xi_0,\xi_1,\xi_2,\xi_3\)，指其在 \(M\) 点的取值，由 §. 166 给为
\[
\tag{46}
 \xi_i=\rho\left(\frac{d\alpha_i}{ds}+\frac{x_i}{R}\right).
\]
切线与主法线确定密切平面；以 \(\lambda_0,\lambda_1,\lambda_2,\lambda_3\) 表示这个平面的坐标，亦即曲线副法线的方向余弦。行列式
\[
\tag{47}
\left|
\begin{matrix}
 x_0&x_1&x_2&x_3\\
 \alpha_0&\alpha_1&\alpha_2&\alpha_3\\
 \xi_0&\xi_1&\xi_2&\xi_3\\
 \lambda_0&\lambda_1&\lambda_2&\lambda_3
\end{matrix}\right|
\]
于是是一个正交代换的行列式；为完全确定这些 \(\lambda\)，约定该行列式的值必须为正，即 \(=+1\)。在 \((x_i)\) 处的密切平面方程为
\[
\left|
\begin{matrix}
 X_0&X_1&X_2&X_3\\
 x_0&x_1&x_2&x_3\\
 \dfrac{dx_0}{ds}&\dfrac{dx_1}{ds}&\dfrac{dx_2}{ds}&\dfrac{dx_3}{ds}\\
 \dfrac{d^2x_0}{ds^2}&\dfrac{d^2x_1}{ds^2}&\dfrac{d^2x_2}{ds^2}&\dfrac{d^2x_3}{ds^2}
\end{matrix}\right|=0,
\]
其中 \(X_i\) 表示动点坐标。读者容易验证，采用 §. 3 对密切平面的定义，或采用该节末尾所述的其他定义之一，也得到同一方程。

密切平面的概念既定，就可完全按照 §. 4 定义第二曲率或挠率 \(\frac1T\)，并推出挠率的值
\[
\tag{48}
 \frac1{T^2}=\mathop{\mathrm{S}}\left(\frac{d\lambda_i}{ds}\right)^2.
\]
随后立即建立 \emph{Frenet 公式}。这些公式把行列式 (47) 各元素对弧长的导数，用这些元素本身以及第一、第二曲率半径 \(\rho,T\) 表出。式 (45)、(46) 已经给出其中两组。现在由恒等式
\[
 \mathop{\mathrm{S}}x\lambda=0,
 \qquad \mathop{\mathrm{S}}\alpha\lambda=0,
 \qquad \mathop{\mathrm{S}}\lambda^2=1,
\]
对 \(s\) 求导并利用式 (45)、(46)，得到
\[
 \mathop{\mathrm{S}}x\frac{d\lambda}{ds}=0,
 \qquad \mathop{\mathrm{S}}\alpha\frac{d\lambda}{ds}=0,
 \qquad \mathop{\mathrm{S}}\lambda\frac{d\lambda}{ds}=0
\]
因而
\[
 \frac{d\lambda_i}{ds}=K\xi_i\qquad (i=0,1,2,3),
\]
其中 \(K\) 是比例因子。

由式 (48) 得 \(K^2=\frac1{T^2}\)，我们取 \(K=\frac1T\)；这样赋予挠率的不仅是绝对值，还有确定的符号。因此
\[
 \frac{d\lambda_i}{ds}=\frac{\xi_i}{T}.
\]
现在只剩下计算导数 \(\frac{d\xi_i}{ds}\)；利用正交行列式 (47) 的性质和已经得到的公式，便可立即完成。

综上，椭圆几何中的 Frenet 公式为
\[
\tag{49}
 \frac{dx_i}{ds}=\frac{\alpha_i}{R},
 \qquad
 \frac{d\alpha_i}{ds}=\frac{\xi_i}{\rho}-\frac{x_i}{R},
 \qquad
 \frac{d\xi_i}{ds}=-\frac{\alpha_i}{\rho}-\frac{\lambda_i}{T},
 \qquad
 \frac{d\lambda_i}{ds}=\frac{\xi_i}{T},
 \qquad i=0,1,2,3.
\]
从这些基本公式出发，完全仿照第一章 §§. 8、9 的步骤，可推出在椭圆空间中，一条曲线的形状也完全由其\emph{内蕴方程}
\[
 \frac1\rho=f(s),
 \qquad \frac1T=\varphi(s),
\]
确定；这两个方程可以任意给定。

采用与上述完全类似的方法，可以建立双曲空间中的 Frenet 公式。沿用同一记号，这些公式写成
\[
\tag{49*}
 \frac{dx_i}{ds}=\frac{\alpha_i}{R},
 \qquad
 \frac{d\alpha_i}{ds}=\frac{\xi_i}{\rho}+\frac{x_i}{R},
 \qquad
 \frac{d\xi_i}{ds}=-\frac{\alpha_i}{\rho}-\frac{\lambda_i}{T},
 \qquad
 \frac{d\lambda_i}{ds}=\frac{\xi_i}{T},
\]
此时 \(x,\alpha,\xi,\lambda\) 之间的关系为
\[
 x_0^2-\sum x_i^2=1,
 \quad \sum\alpha_i^2-\alpha_0^2=1,
 \quad \sum\xi_i^2-\xi_0^2=1,
 \quad \sum\lambda_i^2-\lambda_0^2=1,
\]
\[
 \sum x_i\alpha_i-x_0\alpha_0=0,
 \qquad
 \sum\alpha_i\xi_i-\alpha_0\xi_0=0,
 \quad \text{等等。}
\]
此外，只须在椭圆几何的式 (49) 中把 \(R\) 换成 \(iR\)，把 \(x_1,x_2,x_3\) 换成 \(-ix_1,-ix_2,-ix_3\)，并把 \(\alpha_0,\xi_0,\lambda_0\) 分别换成 \(i\alpha_0,i\xi_0,i\lambda_0\)，便立即得到式 (49*)\srcfn{1}{关于式 (49*) 的直接推导，参见 A. Bazzaroni 的论文 \emph{Le formole del Frenet in geometria iperbolica e loro applicazioni}（Bologna，Gamberini，1897）。}。
